\section{\raggedright Mass readout and consistency constraints}
\paragraph{A necessary norm relation.}
Let $x_i=\sqrt{m_i}\geq0$, with $x\ne0$, and define
\begin{equation}
 K=\frac{\sum_i x_i^2}{(\sum_i x_i)^2},\qquad
 A=\frac{x_0+x_1+x_2}{3},\qquad x=A(1,1,1)+d.
\end{equation}
Since $\sum_i d_i=0$, the constant part and $d$ are orthogonal, so
\begin{equation}
 K=\frac13+\frac{\|d\|^2}{9A^2},\qquad
 K=\frac23\ \Longleftrightarrow\ \|d\|^2=3A^2.
\end{equation}
Thus the Koide value requires equality of the symmetric and traceless
squared norms. Cyclic symmetry preserves these norms separately and
does not fix their ratio. For example, $x=(1,2,3)$ gives $K=7/18$,
while equal roots give $K=1/3$.

\paragraph{A conditional HC representation.}
On Euclidean three-root space set $e=(1,1,1)/\sqrt3$ and let $J_m$
be an orthogonal involution. For $C_m=(I+J_m)/2$, suppose
$x=\alpha C_me$ is nonzero and componentwise nonnegative, with
$\alpha>0$. Writing $h=\langle e,J_me\rangle$ gives
\begin{equation}
 \|C_me\|^2=\frac{1+h}{2},\qquad
 K=\frac{2}{3(1+h)}.
\end{equation}
Consequently $h=0$ is the additional condition for Koide. The projector
axioms permit $J_m=I$, for which $h=1$ and $K=1/3$; they do not select
the desired value. Choosing $J_m$ after specifying the target roots
cannot supply the missing physical selection rule.

An explicit route from the original carrier to this representation
would require a bounded coisometry $B:H\to\mathbb R^3$, $BB^*=I_3$,
whose kernel is $J$-invariant. Then $J_m=BJB^*$ is an orthogonal
involution and $BC=C_mB$. If an independently supplied seed satisfies
$Bf_0=e$ and the physical root law is $x=\alpha BCf_0$, the criterion
becomes
\begin{equation}
 \langle B^*e,JB^*e\rangle=0.
\end{equation}
Neither the name of a readout nor the number of its components supplies
$B$, its metric control, this seed law or the orthogonality selection.
Moreover, a purely odd input obeys $Cf=0$ and cannot yield nonzero roots
through this direct compensated linear route.

\paragraph{Positive nonlinear readout is a separate requirement.}
For two complex amplitudes $\chi=(u,v)$ define
$X_0=|u|^2+|v|^2$, $X_1=|u|^2-|v|^2$,
$X_2=2\operatorname{Re}(\bar uv)$ and
$X_3=2\operatorname{Im}(\bar uv)$. Then
$X_0^2=X_1^2+X_2^2+X_3^2$. With the three-cycle $U$, an illustrative
Hermitian root operator is
\begin{equation}
 L_\eta=X_0\frac{I_3}{\sqrt3}
 +\eta X_1\frac{U+U^*}{\sqrt6}
 +\eta X_2\frac{i(U-U^*)}{\sqrt6},\qquad \eta\geq0.
\end{equation}
Only where $L_\eta\geq0$ may its eigenvalues be read as principal
roots. There
\begin{equation}
 K=\frac{1+\eta^2[1-(X_3/X_0)^2]}3.
\end{equation}
Positivity on every input requires $\eta\leq1/\sqrt2$ and gives
$K\leq1/2$. The isometric coefficient $\eta=1$ can reach $2/3$
on restricted positive phase chambers, but the preparation and
preservation of such a chamber is additional physical information.
Squaring a signed $L_\eta$ does not remove this issue: its physical
principal root becomes $|L_\eta|$, changing the root sum.

The recent HC-balance comparison illustrates the transfer problem
more directly. If the two-mode grading is $J=\operatorname{diag}(1,-1)$,
the balanced state $(1,i)/\sqrt2$ gives $X_1=X_2=0$ and hence $K=1/3$
for this readout. More generally, an antipodal pair of root operators
$cI_3\pm D$, with $c>0$ and $\operatorname{Tr}D=0$, cannot both
be positive Koide readouts. Both positive implies eigenvalues of $D$
in $[-c,c]$ and $\operatorname{Tr}D^2\leq2c^2$, so both ratios
are at most $5/9$. This is a restriction on this isotropic readout
pair, not on arbitrary mass mechanisms. It shows why state balance
and the positive three-generation transfer must be derived together.
\subsection{Why a transition norm is a new readout law}
For a closed seed $P_c\psi=\psi$, the closure commutator satisfies
\begin{equation}
 [\widetilde R,P_c]\psi=(I-P_c)\widetilde R\psi=D\psi.
\end{equation}
The right-hand side lies in the orthogonal complementary space. Hence
$\langle\psi,[\widetilde R,P_c]\psi\rangle=0$, even if its
norm is nonzero. Replacing a commutator expectation by
$\|D\psi\|^2$ therefore changes the observable. It is a legitimate
candidate only when that change is stated and physically justified.

The positive matrix $Q=S^*D^*DS$ measures a tendency to leave the
selected closure space under the supplied operator. It does not
assert that an evolving physical state actually leaves that space.
The operator might be a diagnostic rather than the time generator.
The physical mass law and the preservation of the admitted state
domain must be supplied independently.

The distinction between closure and HC parity gives another necessary
test. Under $[P_c,J]=0$, define
\begin{equation}
 R_o=\tfrac12(\widetilde R-J\widetilde RJ),\qquad
 D_o=\tfrac12(D-JDJ)=(I-P_c)R_oP_c.
\end{equation}
Only $D_o$ changes HC grading. A pacer commuting with $J$ has
$R_o=0$ and cannot supply this channel, even if it has nonzero
closure defect. A selfadjoint comparison coupling can be formed as
$A_o=D_o+D_o^*$, but forming it is an additional operation; its
appearance in the native equations is not established by its definition.

\subsection{Closure penalties do not determine the missing channel}
Let $M$ be a bounded represented matrix and split the carrier into
the range and kernel of $P_c$. Write
\begin{equation}
 M=\begin{pmatrix}A&F\\G&H\end{pmatrix},\quad
 K_c=[M,P_c]=\begin{pmatrix}0&-F\\G&0\end{pmatrix},\quad
 E_c=\begin{pmatrix}0&F\\G&0\end{pmatrix}.
\end{equation}
Direct block multiplication gives
\begin{equation}
 K_c^*K_c=E_c^*E_c=
 \begin{pmatrix}G^*G&0\\0&F^*F\end{pmatrix}.
\end{equation}
Thus a minimal positive penalty
$H_c=\rho K_c^*K_c+\sigma E_c^*E_c$ is block diagonal and
commutes with $P_c$. Its compression can measure the norm of a
transition generated by $M$, but the penalty itself has no off-diagonal
closure channel. Positivity of the penalty cannot replace the missing
represented pacer block.

If the pacer were a norm limit of polynomials in an exactly closed
$M,M^*$, then it would commute with $P_c$ as well and its transition
would vanish. Conversely, exact closure of $M$ says nothing about
an independently specified pacer. One needs a bridge between those
operators before transferring a closure statement from one to the other.

There is a particularly sharp freedom result even when the pacer is
unitary. Given any positive contraction $0\leq Q\leq I_3$, define
\begin{equation}
 R_Q=\begin{pmatrix}\sqrt{I_3-Q}&-\sqrt Q\\
                         \sqrt Q&\sqrt{I_3-Q}\end{pmatrix},\qquad
 P_c=\operatorname{diag}(I_3,0).
\end{equation}
The square roots commute because both are functions of $Q$.
Block multiplication gives $R_Q^*R_Q=I_6$, and the transition from
the first three coordinates to the complement has Gram matrix $Q$.
Therefore unitarity, positive readout and the dimension three leave
every positive contraction spectrum available. These assumptions
do not select a hierarchy or the Koide radius.

\subsection{When a Gram law follows from a gapped block}
A possible physical route can be expressed more concretely at a
frozen background. Assume a canonically normalized Hermitian mass
operator on a retained three-dimensional space and its complement:
\begin{equation}
 \Omega_{\rm tr}=\begin{pmatrix}A&V^*\\V&H_\perp\end{pmatrix},
 \qquad A=A^*,\quad H_\perp=H_\perp^*\geq\Delta I,
 \quad \Delta>0.
\end{equation}
For an eigenvalue $|\lambda|<\Delta$, solve the lower equation to obtain
\begin{equation}
 y=-(H_\perp-\lambda I)^{-1}Vx,
\end{equation}
and substitute into the upper equation:
\begin{equation}
 [A-V^*(H_\perp-\lambda I)^{-1}V]x=\lambda x.
\end{equation}
This is an energy-dependent Schur complement. The resolvent identity
controls the error made by replacing it by its static value:
\begin{equation}
 \|A_{\rm eff}(\lambda)-A_{\rm eff}(0)\|
 \leq\frac{\|V\|^2|\lambda|}{\Delta(\Delta-|\lambda|)}.
\end{equation}
The bound uses the two inverse norms
$\|(H_\perp-\lambda I)^{-1}\|\leq(\Delta-|\lambda|)^{-1}$
and $\|H_\perp^{-1}\|\leq\Delta^{-1}$. It is a controlled
comparison only in the stated gap regime.

With the further assumptions $A=0$, $H_\perp=\Delta I$ and
$V=\mu B$, each eigenvalue $q_i$ of $Q=B^*B$ gives a light
signed eigenvalue $-m_i$, where
\begin{equation}
 m_i=\frac{\sqrt{\Delta^2+4\mu^2q_i}-\Delta}{2}
 =\frac{\mu^2}{\Delta}q_i
  +O\left(\frac{\mu^4q_i^2}{\Delta^3}\right).
\end{equation}
Thus a Gram mass-magnitude law can be a leading approximation in
a canonical Hermitian Dirac-type reading. The signed effective
operator is negative in this comparison; its interpretation as a
positive physical mass magnitude is an additional reading, not
a sign that can be dropped without comment.

The exact eigenstate is dressed. Its complementary norm satisfies
\begin{equation}
 \frac{\|y\|^2}{\|x\|^2}=\frac{m_i}{\Delta+m_i},\qquad
 \|y\|^2=\frac{m_i}{\Delta+2m_i}
\end{equation}
after normalization. Enforcing strict closure $y=0$ would instead
require $Vx=0$. A nonzero light mass produced in this zero-direct-block
mechanism is consequently incompatible with exact closure in the
original retained projector. Approximate closure or a separately
defined dressed spectral projector are possible alternatives, but
they are different state constraints.

\subsection{Cyclic compatibility fixes a form, not its coefficients}
Suppose a unitary carrier operation $W$ and the generation three-cycle
$U$ satisfy
\begin{equation}
 WS=SU,\qquad WP_c=P_cW,\qquad W\widetilde R=\widetilde RW.
\end{equation}
Then $BU=WB$ and hence $QU=UQ$. The positive root $L=\sqrt Q$
also commutes with $U$. In three dimensions the three distinct Fourier
eigenvalues of $U$ imply that every such Hermitian root has the form
\begin{equation}
 L=aI_3+zU+\overline zU^*,\qquad a>0,\quad z=re^{i\theta}.
\end{equation}
Its eigenvalues are
\begin{equation}
 \ell_k=a+2r\cos(\theta+2\pi k/3).
\end{equation}
Trace orthogonality gives
\begin{equation}
 \operatorname{Tr}L=3a,\qquad
 \operatorname{Tr}L^2=3a^2+6r^2,\qquad
 K=\frac13+\frac{2r^2}{3a^2}.
\end{equation}
The Koide value therefore requires $r=a/\sqrt2$. Cyclic compatibility
does not supply this coefficient ratio. Conversely, once the native
$Q$ were represented, its coefficients would be determined without a
mass fit by
\begin{equation}
 a=\tfrac13\operatorname{Tr}\sqrt Q,\qquad
 z=\tfrac13\operatorname{Tr}(U^*\sqrt Q).
\end{equation}
The native burden is thus a joint operator calculation, not the
selection of a familiar trigonometric parametrization.

At the Koide radius positivity gives three phase chambers,
\begin{equation}
 \theta\in\bigcup_{j=0}^2
 [2\pi j/3-\pi/12,\ 2\pi j/3+\pi/12]
 \pmod{2\pi}.
\end{equation}
They follow by intersecting the inequalities
$\cos(\theta+2\pi k/3)\geq-1/\sqrt2$ for all three $k$.
At the boundary $\theta_b=3\pi/4$, write
$\theta=\theta_b-\delta$ inside the chamber. The small root is
$\ell_0=a\delta+O(\delta^2)$, so its candidate mass is
$m_0=\kappa a^2\delta^2+O(\delta^3)$. A light eigenmass can
therefore arise near a boundary in this comparison. Neither the
distance to the boundary nor its dynamical preservation has yet been
selected by the native theory.

\subsection{Deriving the positivity obstruction in the quadratic model}
For $\eta\geq0$, the quadratic root model can be written as a linear function of the
two-mode density matrix $\rho=\chi\chi^*$. In the Fourier basis of
$U$, its signed root eigenvalues are
\begin{equation}
 \ell_k=\frac{X_0}{\sqrt3}
 +\sqrt{\frac23}\eta[X_1\cos\varphi_k-X_2\sin\varphi_k],
 \qquad \varphi_k=2\pi k/3.
\end{equation}
Equivalently $\ell_k=\sqrt3\operatorname{Tr}(F_k\rho)$, where
\begin{equation}
 F_k=\frac{I_2+\sqrt2\eta(\cos\varphi_k\sigma_z-
                               \sin\varphi_k\sigma_x)}{3}.
\end{equation}
The three $F_k$ sum to $I_2$. The Pauli combination has eigenvalues
$\pm1$, so $F_k$ has eigenvalues $(1\pm\sqrt2\eta)/3$.
Positivity for every $\rho\geq0$ is therefore equivalent to
$\eta\leq1/\sqrt2$. Necessity follows by testing a negative
eigenvector if this condition fails; sufficiency follows because
the trace of a product of two positive matrices is nonnegative.
This is the precise sense in which a globally positive linear-density
readout in this family cannot achieve the intended isometric radius.

Replacing the input metric by the pullback metric of this map would
make a chosen nonzero $\eta$ an isometry by definition. It would
not derive its physical value. Nor does covariance under the three-cycle
fix the relative singlet and doublet weights: the constant direction
and the zero-sum plane are separate invariant pieces. A comparison
tensor can assign them different positive weights while preserving
the cyclic symmetry.

The antipodal $5/9$ bound is similarly a property of the actual
root eigenvalues. If $cI_3\pm D$ are both positive, the three
eigenvalues $d_i$ obey $|d_i|\leq c$ and $\sum_i d_i=0$.
The maximum of $\sum_i d_i^2$ on this planar cube section is
$2c^2$, attained at permutations of $(c,0,-c)$. Consequently
\begin{equation}
 \frac{\operatorname{Tr}[(cI_3\pm D)^2]}
      {[\operatorname{Tr}(cI_3\pm D)]^2}
 =\frac{3c^2+\operatorname{Tr}D^2}{9c^2}\leq\frac59.
\end{equation}
Selecting one allowed positive branch can evade the pairwise
requirement. A physical explanation then has to supply that branch
selection rather than infer it from the involution alone.

\subsection{Phase overlap: root averaging and mass averaging differ}
The published overlap comparison provides another explicit test of
what is being averaged. In one common generation basis, let
\begin{equation}
 x_{jk}=a[1+\sqrt2\operatorname{Re}(\omega^ke^{i\delta_j})],
 \quad \omega=e^{2\pi i/3},\quad
 w_j\geq0,\quad\sum_jw_j=1.
\end{equation}
Each preparation is assumed to have nonnegative roots and the same
mean $a$. Averaging the roots before squaring gives
\begin{equation}
 \overline x_k=a[1+\sqrt2\operatorname{Re}(\omega^k Z)],
 \qquad Z=\sum_jw_je^{i\delta_j},\qquad
 K=\frac{1+|Z|^2}{3}.
\end{equation}
For two equally weighted phases $\delta_0\pm h$, this reduces the
radius by $\cos h$ when $|h|<\pi/2$, leaving the central phase.
It follows from a vector average on the unit circle, not from a
new dynamical mass equation.

Averaging squared roots gives a different answer:
\begin{equation}
 \sum_jw_jx_{jk}^2=\overline x_k^2+
                      \sum_jw_j(x_{jk}-\overline x_k)^2.
\end{equation}
The variance term cannot be removed by calling both procedures a
superposition. A coherent coupling of two paths is yet another
specified construction. For $B(\phi)=B_1+e^{i\phi}B_2$,
\begin{equation}
 B(\phi)^*B(\phi)=B_1^*B_1+B_2^*B_2+
 e^{i\phi}B_1^*B_2+e^{-i\phi}B_2^*B_1.
\end{equation}
Here relative phases can affect the operator spectrum through cross
terms. A common phase multiplying all of $B$ cancels from its Gram
matrix. Changing only the prepared probe can change spectral weights
without moving the operator poles. These alternatives must be
distinguished before a transport interpretation is assigned.

An inverse reconstruction can find phase mixtures consistent with
specified masses and an assumed intermediate gap. Its feasibility
region is useful for excluding combinations of assumptions, but the
reconstructed phases and gap are not independent predictions. The
native operator, the admitted preparations and the rule for combining
paths must be fixed before the resulting spectrum is compared with
data as a test of the model.
